% !TeX root = ../../../main.tex
% chapters/sections/chapter5/subsections/two_plane/pd.tex

\subsection{PD(純粋需要)曲線}
\label{sec:chapter5_two_plane_pd}
PD(純粋需要)曲線は需要群の式を統合することで導出される。
PD(純粋需要)曲線は消費者物価指数(CPI)\(p_t^{H\to W}\)を
総消費指数\(c_t^{H\to W}\)の関数として記述したものとなり、
したがって\((c_t^{H\to W},p_t^{H\to W})\)平面上に描かれるCD(消費需要)曲線となる。

\subsubsection{PD(純粋需要)曲線を構成する式}
\label{sec:chapter5_two_plane_pd_components}
PD(純粋需要)曲線を構成する需要群の式は
\ref{sec:chapter5_two_plane_classification_demand}
節で述べられたとおり、以下のようであった。
\begin{itemize}
    \item \textbf{自国コンティンジェント債券\(d_{t+1}^{h\to H}(j')\)のオイラー方程式} \eqref{form:chapter5_two_plane_classification_demand_lambda_H_i_H_eq}
    \item \textbf{所得の限界効用\(\lambda_t^H\)の式} \eqref{form:chapter5_two_plane_classification_demand_lambda_H_p_H_W_c_H_W_eq}
\end{itemize}

\subsubsection{曲線の導出}
\label{sec:chapter5_two_plane_pd_derivation}
自国コンティンジェント債券のオイラー方程式
\eqref{form:chapter3_representative_household_euler_contingent_home_representative_lambda_H_i_H_eq}
は以下のようであった。
\begin{equation}
\label{form:chapter3_representative_household_euler_contingent_home_representative_lambda_H_i_H_eq}
    \lambda_t^H=\beta_t^H(1+i_t^H)\operatorname{E}_t[\lambda_{t+1}^H]
\end{equation}
この左辺の\(\lambda_t^H\)に所得の限界効用\(\lambda_t^H\)の式
\eqref{form:chapter3_representative_household_lambda_home_representative_lambda_H_p_H_W_c_H_W_eq}
\begin{equation}
\label{form:chapter3_representative_household_lambda_home_representative_lambda_H_p_H_W_c_H_W_eq}
    \lambda_t^H=\frac{1}{p_t^{H\to W}c_t^{H\to W}}
\end{equation}
を代入すると以下のようになる。
\begin{equation}
\label{form:chapter5_two_plane_pd_derivation_p_H_W_c_H_W_eq}
    \frac{1}{p_t^{H\to W}c_t^{H\to W}}=\beta_t^H(1+i_t^H)\operatorname{E}_t[\lambda_{t+1}^H]
\end{equation}
これを名目総消費\(p_t^{H\to W}c_t^{H\to W}\)について解くと以下の\textbf{PD(純粋需要)曲線}が得られる。
\begin{equation}
\label{form:chapter5_two_plane_pd_derivation_p_H_W_c_H_W_modified_eq}
    p_t^{H\to W}=\frac{1}{c_t^{H\to W}}\times\frac{1}{\beta_t^H(1+i_t^H)\operatorname{E}_t[\lambda_{t+1}^H]}
\end{equation}
これにより消費者物価指数(CPI)\(p_t^{H\to W}\)と総消費指数\(c_t^{H\to W}\)は反比例の関係にあることがわかる。
したがって\((c_t^{H\to W},p_t^{H\to W})\)平面において
PD(純粋需要)曲線は右下がりの\textbf{直角双曲線}として描かれる。
この曲線の位置は右辺の名目利子率\(i_t^H\)や来期の期待限界効用\(\operatorname{E}_t[\lambda_{t+1}^H]\)によって決定される。
そのため金融政策はこれらの変数を制御することによりこの曲線を移動させ総消費指数\(c_t^{H\to W}\)を調整する。